\section{Physical model and scale analysis}
\label{sec:model}

\subsection{Governing equations}
\label{sec:model-governing}

In the electrode region, the plasma interacts with the relatively cold, absorbing electrode surfaces. This interaction introduces strong gradients in the plasma properties, causing the plasma to progressively depart from LTE.  We consider plasma to be composed of electrons, singly-charged ions, and neutral atoms for the model. The presence of multiply charged species is neglected as justified in \cite{almeida2000simulation}. To estimate the characteristic length scales associated with these departures, we consider the adjacent LTE plasma at two representative temperatures, \(T=10{,}000~\mathrm{K}\) and \(T=20{,}000~\mathrm{K}\), which are typical of the anode and cathode regions, respectively, in a high-intensity atmospheric-pressure arc discharge \cite{Lowke2006,Hsu1983,LiangTrelles}.

In an ionized gas, the electron mean free path can be approximated as \(\lambda_e = 1/(n_e\overline{Q}_{ee}+n_a\overline{Q}_{ea})\), where \(n_e\) and \(n_a\) are the electron and neutral-atom number densities, respectively, and \(\overline{Q}_{ee}\) and \(\overline{Q}_{ea}\) are the corresponding momentum-transfer collision cross sections for electron--electron and electron--atom collisions. Electron--electron collisions are Coulomb collisions, for which the averaged momentum-transfer cross section is given by \(\overline{Q}_{ee}=e^4\ln\Lambda/[32\pi\varepsilon_0^2(k_BT_e)^2]\) \cite{mitchner1973partially}. Here, \(\ln\Lambda\) is the Coulomb logarithm, given by \(\ln\Lambda=1.24\times10^7 T_e^{3/2}n_e^{-1/2}\), where \(T_e\) is in K and \(n_e\) is in \(\mathrm{m}^{-3}\).

The average electron--neutral momentum-transfer cross section is of the order \(0.2\times10^{-19}~\mathrm{m}^2\) \cite{lxcat} for the electron temperatures considered. Using the corresponding LTE electron and neutral densities, the resulting electron mean free path is of order \(\lambda_e\sim10^{-6}~\mathrm{m}\) in the electrode regions. The characteristic length scale for electron--heavy-particle energy relaxation, \(\lambda_T\), is approximated as\(\frac{\lambda_e}{\sqrt{2M/m_e}}\sim2\times10^{-4}~\mathrm{m}\), where \(M\) is the mass of the heavy species (argon in this case) and \(m_e\) is the electron mass. Thus, \(\lambda_T\sim10^{-4}~\mathrm{m}\).Lowke and Tanaka \cite{Lowke2006} identified a characteristic diffusion length scale for the electron and ion densities of \(\lambda_d\sim10^{-4}~\mathrm{m}\) in the electrode regions.

At the smallest spatial scale, the Debye length calculated from the LTE number densities is of order \(\lambda_{De}\sim10^{-8}~\mathrm{m}\), which characterizes the spatial extent over which charge separation and space-charge effects become important. The ion mean free path provides another relevant scale and can be estimated as \(\lambda_i=1/(n_a\overline{Q}_{ia})\), where \(\overline{Q}_{ia}\) is the momentum-transfer cross section for ion--atom collisions. For argon, we use the approximate fit \(\overline{Q}_{ia}=2[7\times10^{-10}-0.6\times10^{-10}\ln(2k_BT_h)]\) \cite{benilov1998ionization}, where \(k_B\) is the Boltzmann constant and \(T_h\) is the heavy-particle temperature. Since the heavy particle temperature vary from LTE condtions $\sim 1000 \mathrm{K}$ to $\sim 1000 \mathrm{K}$, to electrode temp and since we are interested in region close to the electrode, we use a  temperature of \(T_h\sim1000~\mathrm{K}\), which gives an ion mean free path of order \(\lambda_i\sim10^{-7}~\mathrm{m}\).

The resulting hierarchy of characteristic length scales in the near-electrode region is therefore \(\lambda_{De}\ll\lambda_i\ll\lambda_e\ll\lambda_d\sim\lambda_T\). These estimates indicate that the near-electrode region over which significant departures from LTE develop extends over a characteristic length scale of order \(10^{-4}~\mathrm{m}\). Other modeling results results and analysis shows the same \cite{Almeida2008,DinulescuPfender1980,benilov1995model,hsu1983analysis}. Consistent with this estimate, Lowke and Tanaka \cite{Lowke2006} showed that excluding a near-electrode region of approximately \(10^{-4}~\mathrm{m}\) is sufficient to model the remaining plasma using an LTE approximation. This separation of scales provides the basis for treating the bulk arc with an LTE model while explicitly resolving the nonequilibrium physics in the near-electrode region with a 1D model due to the small size of the near-electrode region.



We therefore model the electrode region as a one-dimensional domain, with the electrode surface at one end and the bulk plasma at the other. Within this region, fluid-dynamic effects associated with bulk plasma motion are neglected, since the near-electrode region is typically contained within the hydrodynamic boundary layer adjacent to the electrode. The neutral species are assumed to be stationary and in thermal equilibrium with the electrode surface and with each other. The ions are assumed to remain in thermal equilibrium with the neutrals because their mean free path is small compared with the characteristic length scale of the temperature variation. Accordingly, the heavy-particle temperature is prescribed using an exponential profile that decreases from the bulk plasma temperature to the electrode temperature,

$$
T_h(x)=T_{\mathrm{bulk}}+\left(T_{\mathrm{wall}}-T_{\mathrm{bulk}}\right)\exp\left(-\frac{x}{L_T}\right),
$$

where \(L_T\) is a characteristic thermal length scale. This exact ion temperature profile makes a little difference on the model because the coliision frequency if ion-neutral has a weak dependence on collision energy. The electron number density and temperature, and the ion number density, are therefore solved self-consistently, while the heavy-particle temperature is prescribed. The neutral number density is not solved independently but is obtained from the ideal-gas law using the prescribed heavy-particle temperature and the constant system pressure. The system is closed by Poisson's equation for the electric potential, which is solved throughout the entire computational domain.

In addition, inertial effects are neglected, and the frictional, pressure-gradient, and electric-field forces are assumed to balance locally for each charged species. This approximation corresponds to a quasi-steady drift-diffusion description in which the species flux adjusts rapidly to changes in the local plasma conditions. It is valid when the momentum-relaxation time is sufficiently short compared with the characteristic timescale of the macroscopic evolution, or equivalently when the associated momentum-relaxation length is small compared with the characteristic spatial scale of the plasma properties. The validity of this approximation is discussed in Section~\textcolor{red}{fill later}. The resulting quasi-steady flux is then used in the species continuity equations to determine the spatial evolution of the charged-particle densities.

The continuity equation for species \(s\) is given by
\begin{equation}
\frac{\partial n_s}{\partial t}+\nabla\cdot\boldsymbol{\Gamma}_s=S_s,
\qquad
\boldsymbol{\Gamma}_s=z_s\mu_s n_s\mathbf{E}-D_s\nabla n_s,
\label{eq:continuity}
\end{equation}
where \(z_s=\pm1\) is the signed charge number, \(\mu_s\) is the mobility, \(D_s\) is the diffusion coefficient, and \(S_s\) is the net ionization--recombination source term. For the simple two-body ionization/recombination scheme considered here, the source terms for electrons and ions are equal in magnitude and opposite in sign. The flux \(\boldsymbol{\Gamma}_s\) is the particle flux of species \(s\); its detailed formulation and the calculation of the corresponding transport properties are described in Appendix~\textcolor{red}{fill later}. The ionization and recombination processes are modeled using the formulation of Hoffert and Lien \cite{hoffert1967quasi}.

The electric field is obtained from Poisson's equation,
\begin{equation}
\nabla\cdot(\varepsilon_0\mathbf{E})=e(n_i-n_e),
\label{eq:poisson}
\end{equation}
which is solved over the entire computational domain, thereby allowing the transition from the quasi-neutral bulk plasma to the space-charge region to be captured without imposing a priori sheath boundaries.

The electron energy is governed by a separate equation of the same convection--conduction--Joule-heating--energy-exchange form commonly used in two-temperature bulk-arc models \cite{HsuPfender1983b,Freton2012},
\begin{equation}
\frac{3}{2}\frac{\partial(n_eT_e)}{\partial t}
+\nabla\cdot\left(\frac{5}{2}T_e\boldsymbol{\Gamma}_e-\kappa_e\nabla T_e\right)
={\mathbf{E}\cdot\mathbf{J}}-\dot{E}_{eh}-\dot{E}_{\rm inel},
\label{eq:eenergy}
\end{equation}
where the terms on the right-hand side represent Joule heating, elastic electron--heavy-particle energy exchange, and inelastic energy losses, respectively. The elastic energy-exchange term is given by \(\dot{E}_{eh}=n_e\delta\nu_{en}\tfrac{3}{2}k_B(T_e-T_{\rm gas})\), where \(\delta=2m_e/M\), while \(\dot{E}_{\rm inel}\) accounts for energy losses associated with inelastic processes, including ionization. Radiation losses are usually small \cite{DinulescuPfender1980,Almeida2008} and thus not included.




Heavy-particle temperature is either taken as a fixed background gas
temperature $T_{\rm gas}(x)$ (the regime in which the present
implementation, Section~\ref{sec:methods}, currently operates) or evolved by
its own energy equation as in \cite{HsuPfender1983b,Almeida2008} when a
fully coupled two-temperature heavy-species treatment is required.


%%%%%%%%%%%%%%%%%%%%%%%%%%%%%%%%%%%%%%%%%%%%%%%%%%%%%%%%%%%%%%%%%%%%%%%%%%%%%%%%%%%
\section{Appendix A. Transport Coefficients}

The electron mobility is obtained from a two-term Boltzmann solution of the BOLSIG+-type, assuming a Maxwellian electron energy distribution, and is tabulated in reduced form as \(\mu_eN(T_e)\) and \(D_eN(T_e)\), rather than being evaluated from a closed-form kinetic-theory expression. This treatment is appropriate in the weakly ionized regime, where electron--neutral collisions dominate and the electron energy distribution is close to Maxwellian. The validity of the Maxwellian assumption is discussed in Section~\textcolor{red}{fill later}. However, in the core region of the arc, particularly near the cathode, electron--ion Coulomb collisions can become significant. To provide a continuous description between these regimes, an electron--ion Coulomb correction is incorporated into the BOLSIG+-based transport coefficients.

The electron--neutral collision frequency is obtained from the BOLSIG+-based mobility as \(\nu_{en}=e/(m_e\mu_{en})\), while the electron--ion collision frequency is calculated using the NRL Plasma Formulary as \(\nu_{ei}=2.91\times10^{-12}n_e[\mathrm{m}^{-3}]\ln\Lambda\,T_e[\mathrm{eV}]^{-3/2}~\mathrm{s}^{-1}\). In the weakly ionized limit, where electron--neutral collisions dominate, the electrical conductivity is \(\sigma_{\rm en}=n_ee^2/(m_e\nu_{en})\), whereas for a fully ionized plasma the electron--ion contribution can be written as \(\sigma_{\rm ei}=1.975\,n_ee^2/(m_e\nu_{ei})\) \cite{mitchner1973partially}. Combining these expressions gives the effective electron mobility
\begin{equation}
\mu_e=\frac{e}{m_e\left(\nu_{en}+\nu_{ei}/1.975\right)}.
\end{equation}
Thus, the BOLSIG+-based electron mobility is corrected by the factor \(\nu_{en}/(\nu_{en}+\nu_{ei}/1.975)\), yielding the effective mobility used in the model. The same correction factor is applied to the electron diffusion coefficient, \(D_e\), to account for the additional momentum loss associated with electron--ion Coulomb collisions. The resulting electrical conductivity is compared with the LTE conductivity in Fig.~\textcolor{red}{fill later}.

The ion--neutral momentum-transfer cross section is obtained from the fitted expression
\(\overline{Q}_{ia}=2[7\times10^{-10}-0.6\times10^{-10}\ln(2k_BT_h)]\) \cite{benilov1998ionization}. The low-field ion mobility is then calculated using the Mason--Schamp equation \cite{mason1958mobility},
\begin{equation}
\mu_i =
\frac{3e}{16N}
\sqrt{\frac{2\pi}{\mu k_B T_{\mathrm{gas}}}}
\frac{1}{\Omega^{(1,1)}(T_{\mathrm{gas}})},
\end{equation}
where \(N\) is the neutral number density, \(\mu\) is the reduced mass of the ion--neutral pair, and \(\Omega^{(1,1)}\) is the corresponding collision integral. This expression gives the ion mobility in the low-field limit, where the ion drift velocity is proportional to the applied electric field.

To account for the field dependence of the ion mobility, a closed-form two-temperature Wannier formulation is employed. The drift velocity is expressed as \(v_d=\mu(T_{\rm eff})E\), where

$$
\mu(T_{\rm eff})=\mu_0\sqrt{\frac{T_{\rm gas}}{T_{\rm eff}}},
\qquad
k_BT_{\rm eff}=k_BT_{\rm gas}+\frac{1}{3}Mv_d^2,
$$

with \(\mu_0\) denoting the low-field mobility and \(M\) the ion mass. Substitution of these relations yields a quadratic equation in \(v_d^2\), which is solved to obtain the field-dependent drift velocity. In the limit \(E\rightarrow0\), the effective temperature approaches the gas temperature, \(T_{\rm eff}\rightarrow T_{\rm gas}\), and the formulation recovers the low-field mobility, \(\mu\rightarrow\mu_0\).




\subsection{Quasi-Maxwellian electrons near the electrode}
\label{sec:model-scaling}

Even where the heavy species are strongly out of equilibrium, the electron
population can remain closely Maxwellian provided electron--electron and
electron--ion collisions are frequent enough to redistribute directed
(drift) energy into the isotropic thermal spread faster than electron--heavy
collisions or the imposed field can distort it. Following Dinulescu and
Pfender's analysis of the anode boundary layer \cite{DinulescuPfender1980},
two conditions must hold. First, electron--electron/ion thermalization must
dominate electron--neutral momentum transfer, weighted by the mass ratio
that makes each electron--neutral collision inefficient at removing energy:
\begin{equation}
\sum_h \frac{m_e}{m_h}\,\frac{\bar\nu_{eh}}{\bar\nu_{ee}} \ll 1 .
\label{eq:maxw1}
\end{equation}
Second, the ordered drift energy imparted by the field must remain small
compared with the electron thermal energy, weighted by the same mass ratio:
\begin{equation}
\left(\frac{U_e}{\bar c_e}\right)^2 \ll \frac{m_e}{m_h} .
\label{eq:maxw2}
\end{equation}
For a representative near-anode condition, $T_e\sim1$~eV, $n_e\sim
5\times10^{22}\,\mathrm{m^{-3}}$, neutral density $N\sim6\times10^{23}\,
\mathrm{m^{-3}}$, electron--atom momentum-transfer cross section
$\bar Q_{ea}=2\times10^{-20}\,\mathrm{m^2}$, and mean thermal speed
$\bar c_e = 7\times10^5\,\mathrm{m/s}$: the electron--ion/electron--electron
frequency from the NRL formulary (with $\ln\Lambda\approx3$) is
$\nu_{ee}\sim\nu_{ei}\sim5\times10^{11}\,\mathrm{s^{-1}}$, while the
electron--neutral frequency is $\nu_{en} = \bar c_e N \bar Q_{ea} \approx
8\times10^{9}\,\mathrm{s^{-1}}$, i.e.\ $\nu_{en}/\nu_{ee}\sim0.016$; with
$m_e/M_{\rm Ar}\approx1.4\times10^{-5}$, condition~(\ref{eq:maxw1}) gives
$\sim2\times10^{-7}\ll1$. For condition~(\ref{eq:maxw2}), the electrical
conductivity of the layer, $\sigma\approx5000$~S/m, corresponds to an
electron mobility $\mu_e = \sigma/(n_e e)\approx0.6\,\mathrm{m^2/(V\,s)}$;
with a representative current density $J\sim10^7\,\mathrm{A/m^2}$, the
electron drift speed is $U_e = J/(en_e) \approx1.3\times10^3\,\mathrm{m/s}$,
so $(U_e/\bar c_e)^2\approx3\times10^{-6}$, again well below
$m_e/M_{\rm Ar}\approx1.4\times10^{-5}$. Both conditions are comfortably
satisfied, so the electron population is treated as quasi-Maxwellian
throughout the near-electrode layer, and a single electron temperature
$T_e$ (rather than a fully kinetic electron energy distribution) is
retained in equation~(\ref{eq:eenergy}) there as well as in the bulk.

\subsection{Length scales and the thickness of the near-electrode layer}
\label{sec:model-lengths}

The same parameters fix a hierarchy of length scales that separates the
space-charge sheath from the broader near-electrode non-equilibrium region
that must be resolved around it. The electron Debye length,
\begin{equation}
\lambda_{De} = \sqrt{\frac{\varepsilon_0 T_e}{n_e e}} \approx 3\times10^{-8}
\ \mathrm{m} ,
\label{eq:debye}
\end{equation}
sets the scale of the space-charge sheath itself. The electron mean free
path against the combined electron--neutral and electron--electron/ion
collision frequency used above,
\begin{equation}
\lambda_e = \frac{\bar c_e}{\nu_{en}+\nu_{ee}} \approx 1\times10^{-6}\
\mathrm{m} ,
\label{eq:mfp}
\end{equation}
is $\sim30$ times $\lambda_{De}$, confirming that the sheath is thin
compared with a mean free path and can support the non-collisional
electron kinetics discussed in Section~\ref{sec:model-collisionality}. The
electron--heavy energy-relaxation length, over which an electron transfers
its excess thermal energy to the heavy gas through the inefficient
mass-ratio-weighted elastic channel in equation~(\ref{eq:eenergy}), is
larger still,
\begin{equation}
\lambda_T = \frac{\lambda_e}{\sqrt{2m_e/M}} \approx 2\times10^{-4}\
\mathrm{m} .
\label{eq:thermal-length}
\end{equation}
Finally, the diffusion length that controls how far the electron/ion density
profile itself can depart from its local equilibrium value under a given
field, used directly by Lowke and Tanaka's mesh-sizing rule
\cite{Lowke2006}, is $\Lambda = (kT_e/e)/E$; with $T_e\approx1$~V and a
field $E = J/\sigma \approx (10^7\,\mathrm{A/m^2})/(5000\,\mathrm{S/m}) =
2000$~V/m consistent with the parameters above,
\begin{equation}
\Lambda \approx 5\times10^{-4}\ \mathrm{m} .
\label{eq:diffusion-length}
\end{equation}
The four scales satisfy $\lambda_{De} \ll \lambda_e \ll \Lambda \sim
\lambda_T$, the same ordering (Debye length $\ll$ mean free path $\ll$
boundary-layer thickness) reported directly in Table~I of
\cite{DinulescuPfender1980} for the anode boundary layer of a comparable
arc, $\lambda_{De}\sim(2$--$3)\times10^{-8}$~m,
$\lambda_e\sim(1$--$2)\times10^{-6}$~m, boundary-layer thickness
$L\sim0.2$--$0.6$~mm. On this basis we take a near-electrode layer of
thickness $\delta \approx 0.2$~mm --- comparable to both $\Lambda$ and
$\lambda_T$, and consistent with the range of $L$ found in
\cite{DinulescuPfender1980} --- as the region over which the unified,
non-equilibrium model of Section~\ref{sec:methods} must be applied; outside
it, the plasma is close enough to collision-dominated equilibrium for the
LTE bulk treatment of Section  to hold. Note that this
$\sim$0.2~mm layer is not itself the space-charge sheath: it is the outer,
thermal/diffusion-dominated non-equilibrium region (corresponding to layers
2--3 of the four-sublayer structure found by \cite{Almeida2008}) inside
which a much thinner space-charge sheath, of order $\lambda_{De}$ to a few
times $\lambda_{De}$, is nested.

\subsection{Is the sheath collisionless?}
\label{sec:model-collisionality}

A collisionless sheath is the standard simplifying assumption used to close
the NLTE-sheath and two-temperature ionization-layer-sheath approaches of
Section~\ref{sec:intro-electrode}, and it is tempting to expect it to hold
automatically at atmospheric pressure, where collision rates are high.
Almeida, Benilov and Naidis show this expectation is wrong: computing the
number of ion collisions actually incurred while crossing the sheath gives
values from order unity to $10^4$, not $\ll1$, even for atmospheric-pressure
argon \cite{Almeida2008}. The reason is that the sheath thickness itself
scales with the \emph{local} Debye length, $\lambda_{De}\propto
n_e^{-1/2}$, which is set by the arc's own high electron density and
shrinks faster, as the arc becomes hotter and more strongly ionized, than
the neutral mean free paths shrink with pressure (the Ramsauer minimum in
the argon electron--neutral cross section keeps $\lambda_e$ from falling as
quickly as $\lambda_{De}$). A collisionless sheath at high pressure is
therefore not evidence of a crude or incomplete model; it is close to the
generic outcome for a dense atmospheric-pressure arc, and the question that
actually needs answering species by species and quantity by quantity is
\emph{how} collisionless, and for what.

We make this quantitative with four Knudsen numbers, each the ratio of the
relevant mean free path to the resolved sheath width $s$: $\mathrm{Kn}_e =
\lambda_e/s$ (do electrons collide crossing the sheath at all);
$\mathrm{Kn}_\varepsilon = \lambda_T/s$ (electron--heavy energy exchange:
does a local electron energy balance exist inside the sheath);
$\mathrm{Kn}_{ee} = \lambda_{ee}/s$ (electron--electron thermalization:
does a well-defined $T_e$, i.e.\ a locally Maxwellian electron energy
distribution, exist inside the sheath at all -- a distinct and much more
restrictive question than $\mathrm{Kn}_\varepsilon$, since $\lambda_{ee}$
and $\lambda_T$ differ by roughly $10^4$); and $\mathrm{Kn}_i =
\lambda_i/s$ for the ion charge-exchange mean free path. For the converged,
resolved-sheath operating point of the \textsc{PlasmeX} test case discussed
in Section~\ref{sec:methods} ($s = 1.65\times10^{-7}$~m $= 5.6\,
\lambda_{De}$, $T_e\approx0.98$~eV at the wall), these evaluate to
$\mathrm{Kn}_e\approx63$, $\mathrm{Kn}_\varepsilon\approx1.2\times10^4$,
$\mathrm{Kn}_{ee}\approx5.3$, and $\mathrm{Kn}_i\approx0.8$--$1.7$: the
sheath is collisionless for electron momentum and electron--heavy energy
exchange by a wide margin, only marginally collisionless for electron
thermalization, and carries of order one ion collision, consistent with the
lower end of the $10^2$--$10^4$ range quoted by \cite{Almeida2008} for a
somewhat different case. Despite $\mathrm{Kn}_{ee}\sim5$ being only
marginal, an electron population entering a purely retarding, elastic
potential barrier is mapped onto itself by Liouville's theorem: reflection
by an energy-conserving barrier preserves a Maxwellian at unchanged $T_e$
while depleting the density by the Boltzmann factor, giving a testable,
parameter-free reference prediction, $T_e(x) = T_{e,\rm edge} = {\rm
const}$ across the sheath, against which the local closure used in
equation~(\ref{eq:eenergy}) can be checked directly (Section~\ref{sec:methods}).
No such collisionless reference state exists for the ions, which do not see
a retarding barrier and for which $\mathrm{Kn}_i = O(1)$ rules out both the
collisionless limit and full collisional equilibrium; ion behaviour in the
sheath is therefore treated as an explicit target for the resolved
drift-diffusion solution rather than assumed a priori in either limit.
